%% 自動生成（gen.py）
\bktocchap{6-0}{「次の瞬間を決めるルール」で未来を予想する}{天気予報もロケットも、同じ考え方で計算している}{1}{2}{2}{8}{J}
\bktocpart{第1部}{微分方程式の入口}{24}
\bktocchap{6-1}{微分方程式ってなに？}{答えが数ではなく「関数」になる方程式}{1}{2}{2}{26}{C}
\bktocchap{6-2}{表計算で解こう：オイラー法}{80 ℃のお湯は、10分後に何℃？}{1}{2}{2}{66}{C}
\bktocchap{6-3}{積分するだけで解ける形}{落ちるボールの位置を、式で求める}{2}{3}{3}{102}{C}
\bktocchap{6-4}{自分に比例してへる：冷めるお湯}{熱いほど速く冷める。答えは指数関数}{3}{3}{3}{132}{C}
\bktocchap{6-5}{空気抵抗と終端速度}{雨つぶは、なぜどこまでも速くならない？}{1}{2}{3}{176}{C}
\bktocchap{6-6}{ばねと振り子}{小さなゆれなら、周期はゆれの大きさによらない}{4}{3}{3}{206}{C}
\bktocchap{6-7}{ゆれながら止まる：減衰振動}{針がすっと止まるメーターのひみつ}{4}{3}{4}{246}{C}
\bktocpart{第2部}{複素数と双曲線関数}{278}
\bktocchap{6-8}{2乗してマイナスになる数：虚数}{「2回かけて−1」になる数を、回転で考える}{4}{4}{4}{280}{C}
\bktocchap{6-9}{複素平面：かけ算は回転と拡大}{かけ算すると、矢印がくるっと回る}{4}{4}{4}{312}{C}
\bktocchap{6-10}{オイラーの公式}{指数関数と三角関数が、1本の式でつながる}{4}{4}{5}{344}{C}
\bktocchap{6-11}{双曲線関数 cosh・sinh}{たるませた電線の形は、放物線ではない}{3}{4}{5}{386}{C}
\bktocpart{第3部}{式で解く}{418}
\bktocchap{6-12}{減衰振動を式で解く：特性方程式}{2次方程式の判別式が、ゆれ方を決める}{4}{5}{5}{420}{C}
\bktocchap{6-13}{ゆらし続けると：強制振動と共振}{小さな力でも、タイミングが合うと大きくゆれる}{4}{5}{5}{458}{C}
\bktocchap{6-14}{微分方程式がいくつも：連立微分方程式}{つながった2つの水そうの水位は、どう変わる？}{5}{5}{5}{492}{C}
\bktocchap{6-15}{まとめ：物理は「微分方程式＋はじめの状態」で書ける}{ばねも電気も熱も、同じ形の式で書ける}{1}{2}{3}{534}{C}
\par\vspace{2mm}
\bktocbacklist{\bkbi{こたえ}{560}\bkbi{\R{記号}{き|ごう}カード}{575}\bkbi{よく\R{使}{つか}う\R{定数}{てい|すう}・\R{単位}{たん|い}・\R{公式}{こう|しき}}{580}\bkbi{できるようになったことチェック\R{表}{ひょう}}{592}\bkbi{つぎの\R{巻}{かん}へ}{594}\bkbi{さくいん}{597}}
